\documentclass[UTF8,zihao=-4]{ctexart}
\usepackage[a4paper,margin=2.2cm]{geometry}
\usepackage{fontspec}
\usepackage{tabularx}
\usepackage{array}
\usepackage{ragged2e}
\usepackage{fancyhdr}
\usepackage{hyperref}
\setCJKmainfont{Noto Serif SC}
\setCJKsansfont{Noto Sans SC}
\setmainfont{Noto Serif SC}
\setlength{\parindent}{0pt}
\setlength{\parskip}{0.45em}
\pagestyle{fancy}
\fancyhf{}
\fancyhead[L]{商务智能}
\fancyfoot[C]{\thepage}
\hypersetup{colorlinks=true,linkcolor=black,urlcolor=blue}
\begin{document}
\title{商务智能知识点汇总}
\author{}
\date{}
\maketitle
\tableofcontents
\newpage
\section*{1、BI的目标}
\addcontentsline{toc}{section}{1、BI的目标}

BI 的目标主要有两个：一是支持决策，二是改善信息访问。它通过对企业内外部数据的收集、集成、分析和展示，使数据转化为有意义的信息和知识，从而帮助管理者更好地理解企业真实状况，提高决策质量。\par

\newpage
\section*{2、商务智能的定义}
\addcontentsline{toc}{section}{2、商务智能的定义}

商务智能是企业利用现代信息技术，对结构化和非结构化商务数据进行收集、管理、分析和展示，将数据转化为信息、知识和见解，从而改善信息访问，支持管理决策，提升业务流程、绩效和综合竞争力的能力。\par

\newpage
\section*{3、ETL的作用}
\addcontentsline{toc}{section}{3、ETL的作用}

ETL是抽取、转换、装载\par

1. 辨识与主题相关的原始数据\par
也就是从多个数据源中找出对分析主题有用的数据。\par

2. 开发数据抽取策略\par
解决从哪里抽、抽哪些、什么时候抽、怎么抽的问题。\par

3. 保证数据正确和完整\par
数据进入仓库前，需要处理错误、缺失、不一致等问题。\par

4. 将原始数据转换为目标规格\par
包括格式转换、编码转换、单位统一、字段重命名、关键字转换、数据清洗、数据合并、数据汇总等。\par

5. 将数据加载到预定目标区域\par
目标可以包括原子层和集成数据、数据集市、ODS、缓冲区等。\par

ETL 的作用是把来自不同数据源的原始数据抽取出来，经过清洗、转换、集成、汇总等处理，使其变成正确、完整、一致、符合目标规格的数据，并装载到数据仓库、数据集市、ODS 或缓冲区中，为后续 OLAP、报表查询和数据挖掘提供可靠的数据基础。\par

\newpage
\section*{4、数据仓库的作用}
\addcontentsline{toc}{section}{4、数据仓库的作用}

数据仓库是一个面向主题的、集成的、非易失的、时变的数据集合，用于支持经营管理过程中的决策制定。\par

数据仓库的作用可以从几个方面理解：\par

1. 为 BI 提供统一、集成的数据基础\par

企业数据往往来自多个系统、多个部门、多个平台，格式也可能不同。数据仓库把这些分散数据集成起来，形成一个统一的数据环境。\par

2. 支持分析型处理和决策支持\par

数据仓库不是为了替代传统数据库，而是为了适应市场竞争加剧后企业对分析型处理的需要。数据仓库技术正在成为企业信息集成和辅助决策应用的关键技术之一。\par

3. 把分析型环境和事务型环境分离\par

为了提高分析处理和决策支持的效率与有效性，必须把分析型处理及其所需的综合性分析数据，从传统事务型处理和细节性操作数据中分离出来，并按分析型处理需要重新组织，建立单独的分析处理环境。数据仓库正是为建立这种新环境而出现的数据存储和组织技术。\par

4. 支持数据访问、OLAP 和数据挖掘\par

数据仓库还需要支持多种访问方式和分析展示工具，包括查询与报表、桌面 OLAP、关系 OLAP、多维 OLAP、数据挖掘、仪表盘和客户决策支持界面等。\par

数据仓库的作用是为商务智能提供统一、集成、稳定、面向主题、具有历史时间跨度的数据基础。它通过对多源数据进行集成、清洗、聚集和预计算，把事务型数据转化为适合分析型处理的数据环境，从而支持 OLAP、报表查询、数据挖掘和经营管理决策。\par

\newpage
\section*{5、为什么事务处理环境中不适宜建立 DSS 等分析型应用}
\addcontentsline{toc}{section}{5、为什么事务处理环境中不适宜建立 DSS 等分析型应用}

在传统的以数据库为核心的事务处理环境中，不适宜直接建立 DSS 等分析型应用，主要原因有六点。第一，事务处理和分析处理的性能特性不同。事务处理操作时间短、访问数据量小、频率高、并发大，而分析处理运行时间长、访问数据量大、占用资源多，会影响事务系统性能。第二，存在数据集成问题。DSS 需要企业内外部全面、正确的数据，而事务系统通常按部门和应用分散建设，容易产生数据不一致、同名异义、同义异名、单位不一致和“蜘蛛网”等问题。第三，存在动态集成问题。分析数据需要随数据源变化进行周期性刷新，而传统事务处理环境不具备良好的动态集成能力。第四，存在历史数据问题。事务处理通常只保存当前或短期数据，而决策分析需要大量历史数据。第五，存在数据综合问题。事务系统保存大量细节数据，而分析处理需要综合性、汇总性数据，传统事务系统缺乏这种综合能力。第六，存在数据访问问题。事务处理以实时读写和更新为主，分析处理以读操作和定时刷新为主，访问方式不同。因此，应将分析型处理及其所需数据从事务型处理环境中分离出来，重新组织成独立的数据仓库环境，以支持 DSS、OLAP 和数据挖掘\par

\newpage
\section*{6、操作型/分析型数据对比}
\addcontentsline{toc}{section}{6、操作型/分析型数据对比}

\begin{tabularx}{\linewidth}{|*{3}{>{\RaggedRight\arraybackslash}X|}}
\hline
特性 & 操作型数据（DB） & 分析型数据（DW） \\
\hline
定位 & 面向应用的事务处理 & 面向主题的数据分析 \\
\hline
DB设计 & E-R模型 & 星型/雪花模型，数据立方体 \\
\hline
数据 & 当前的、最新的 & 历史的，具有时间跨度 \\
\hline
汇总 & 原始的，细节的 & 集成的，一致的 \\
\hline
视图 & 详细的，关系的 & 总体的，多维的 \\
\hline
操作类型 & 读/写（可变的） & 读（稳定的） \\
\hline
存取请求 & 可预知的 & 事先未知的 \\
\hline
访问记录 & 一次操作少量记录 & 一次操作大量记录 \\
\hline
DB规模 & 100MB \textasciitilde{} GB & TB \\
\hline
工作单位 & 短的，简单事务 & 复杂查询 \\
\hline
性能要求 & 对性能要求高 & 对性能要求较宽松 \\
\hline
\end{tabularx}

\newpage
\section*{7、数据仓库的特征}
\addcontentsline{toc}{section}{7、数据仓库的特征}

数据仓库的四大特征就是：面向主题、集成、非易失、时变。\par

1.1 面向主题\par

数据仓库不是按照具体业务应用或部门系统来组织数据，而是按照管理决策中的分析主题来组织数据。\par

例如主题可以是：\par

商品\par
供应商\par
顾客\par
CRM 中的客户关系管理\par
ERP 中的销售管理、产品质量控制、库存管理等\par


1.2 集成\par

集成是指数据仓库要把来自不同部门、不同系统、不同平台、不同格式的数据统一起来。\par

操作型系统中的数据可能存在：\par

同名异义\par
同义异名\par
单位不一致\par
编码不一致\par
数据类型不一致\par
数据分散在不同系统中\par

数据仓库通过 ETL 把它们清洗、转换、统一后形成一致的数据环境。\par

1.3 非易失 / 稳定\par

数据仓库中的数据主要用于分析处理，通常不是实时修改，而是以查询、分析为主。PPT 中也把分析型数据的操作类型概括为“读，稳定的”。\par


1.4 时变\par

数据仓库保存历史数据，具有时间跨度。操作型数据库主要保存当前数据，而数据仓库关注历史变化，用于研究趋势。PPT 中指出，分析处理更看重 5 到 10 年的历史数据，而事务处理一般只需要当前或短期数据。\par

回答：\par
1.面向主题,指数据仓库中的数据按照较高层次的分析对象进行组织，而不是按照具体应用系统组织。主题是企业信息系统中数据的综合、归类和分析利用对象。\par
2.集成，集成是指数据仓库要把来自不同部门、不同系统、不同平台、不同格式的数据统一起来。数据仓库通过 ETL 把它们清洗、转换、统一后形成一致的数据环境。\par
3.非易失，非易失性是指数据进入数据仓库后一般不再像操作型数据库那样频繁更新、删除或修改，而是相对稳定地保存，用于查询和分析。\par
4.时变，时变性是指数据仓库中的数据反映一段时间内的数据变化，通常带有时间属性，可以支持趋势分析和历史分析。\par

\newpage
\section*{8、数据目标}
\addcontentsline{toc}{section}{8、数据目标}

数据目标：原子层和集成数据、数据集市、操作数据存储 ODS、缓冲区\par

1. 原子层和集成数据\par

原子层可以理解为数据仓库中最基础、最核心的数据存储层。\par

它保存的是经过 ETL 处理后的、较低粒度的数据。所谓“原子”，强调的是数据尽可能接近细节层次，方便后续进行不同层次的综合、汇总和分析。\par

它的作用是：\par

保存历史数据；\par
保存集成后的企业级数据；\par
为上层数据集市、OLAP、报表、数据挖掘提供基础；\par
支持较细粒度的查询和重新汇总。\par

这一层偏向“仓库底座”。如果没有原子层，后续的数据集市、统计汇总和多维分析就缺少可靠基础。\par

2. 数据集市 Data Mart\par

数据集市是面向特定部门、特定用户群、特定商务单元或特定分析应用的数据集合。\par

全局数据仓库通常太大，实际应用中常常按部门或个人建立反映子主题与区域的局部性数据组织，这就是数据集市；例如商品销售数据仓库中可以建立商品采购数据集市、库房使用数据集市、商品销售数据集市等。\par

它的作用是：\par

减轻直接访问企业级数据仓库的压力；\par
为部门或具体主题提供更直接的数据；\par
提高查询和分析效率；\par
面向终端用户的具体分析需求。\par

数据集市可以看作是从数据仓库中再抽取、再组织出来的局部数据环境。数据仓库与数据集市的关系类似于关系数据库中“基表与视图”的关系，数据集市的数据来自数据仓库，是数据仓库中数据的一部分与局部。\par

3. ODS 操作数据存储\par

ODS 全称是 Operational Data Storage，操作数据存储。\par

ODS 是在企业范围内，针对特定主题区域，用于支持战术决策支持的综合数据的更新集合。它具有以下特点：面向特定分析应用、完整性、当前有效、可变的、详尽的。\par

ODS 和数据仓库的区别可以这样理解：\par

数据仓库更强调历史性、稳定性、长期趋势分析；\par
ODS 更强调当前性、可变性、战术决策支持；\par
ODS 中的数据比操作型数据库更集成，但比数据仓库更接近当前业务状态。\par

所以 ODS 常常处在 DB 与 DW 之间，既不像 OLTP 那样只服务日常事务，也不像 DW 那样主要保存长期稳定历史数据。\par

可背诵：\par
ODS 是面向企业范围内特定主题区域、支持战术决策的综合数据更新集合，具有当前有效、可变、详尽、完整和面向特定分析应用等特点。\par

4. 缓冲区 Staging Area\par

缓冲区是 ETL 过程中的临时中间区域。\par

缓冲区是“数据流的中间站”，是“白板式”的，没有特定结构，系统中可能存在多处缓冲区。\par

它的作用主要是：\par

暂存从数据源抽取出来的原始数据；\par
为清洗、转换、合并、去重、格式转换提供操作空间；\par
降低 ETL 过程对源系统和目标仓库的直接影响；\par
支持分阶段处理复杂数据。\par

缓冲区一般不是最终分析用户直接访问的地方，而是 ETL 内部使用的“加工区”。\par

答案：\par
1、原子层和集成数据是数据仓库的基础数据层，保存经过抽取、清洗、转换和集成后的较低粒度历史数据，是上层分析应用和数据集市的数据来源。\par
2、数据集市是面向特定商务单元、部门、应用或用户群的数据组织，是数据仓库中数据的局部和再组织，通常用于满足特定主题的快速分析需求。\par
3、操作数据存储是面向企业范围内特定主题区域、支持战术决策的综合数据更新集合，具有当前有效、可变、详尽、完整和面向特定分析应用等特点。\par
4、缓冲区是数据流进入数据仓库过程中的中间站，主要用于暂存和处理抽取数据，支持数据清洗、转换和集成。它通常没有固定结构，系统中可以存在多个缓冲区。\par

\newpage
\section*{9、数据刷新的四种方法}
\addcontentsline{toc}{section}{9、数据刷新的四种方法}

时间戳、DELTA 文件、建立映象文件、日志文件\par

数据刷新主要有四种方法：时间戳、DELTA 文件、建立映象文件和日志文件。时间戳依赖记录中的修改时间；DELTA 文件记录应用修改操作；映象文件通过两次快照比较差异；日志文件利用 OLTP 数据库日志进行增量刷新。四种方法可以在同一数据仓库系统中结合使用，以适应不同数据源。\par

\newpage
\section*{10、建立数据集市原因}
\addcontentsline{toc}{section}{10、建立数据集市原因}

企业级数据仓库通常规模较大，虽然提供统一数据模式、统一数据表示和统一访问接口，但终端用户在具体部门或应用场景下仍然需要面向局部主题的快速分析环境。网络环境中即使存在多个数据源，最终用户仍可能得不到可用信息。早期数据仓库主要强调多数据源集成，为用户访问多个数据源提供统一视图。因此，可以按照部门、个人、子主题或业务区域建立数据集市，以满足特定商务单元、商务程序或商务应用的需求。\par

\newpage
\section*{11、数据仓库和数据集市的关系}
\addcontentsline{toc}{section}{11、数据仓库和数据集市的关系}

数据仓库与数据集市的关系类似于传统关系数据库系统中“基表”和“视图”的关系。\par

四种数据仓库与数据集市的结构关系：\par

自顶向下的结构\par
自底向上的结构\par
总线结构的数据集市\par
企业级数据集市结构\par

1. 自顶向下结构先构建企业数据仓库，再基于企业数据仓库构建数据集市，优点是数据整合统一，缺点是成本高、见效慢\par

操作型数据 / 外部数据\par
\hspace*{1em}↓\par
企业数据仓库 Enterprise Warehouse\par
\hspace*{1em}↓\par
局部数据集市 Local Data Mart\par
\hspace*{1em}↓\par
部门或应用分析\par

2. 自底向上结构先构建局部数据集市，再合并为企业数据仓库，优点是见效快、启动资金少，缺点是容易造成蜘蛛网和数据不一致\par

局部操作型数据\par
\hspace*{1em}↓\par
局部数据集市 Local Data Mart\par
\hspace*{1em}↓\par
多个数据集市合并\par
\hspace*{1em}↓\par
企业数据仓库 Enterprise Warehouse\par

3. 总线结构不先建立中心数据仓库，而是通过共享维表和事实表把多个数据集市连接起来，优点是支持共享和增量建设，缺点是主要适用于 OLAP 且结构稳定性较差\par

数据源\par
\hspace*{1em}↓        ↓        ↓\par
数据集市A — 共享维表/事实表 — 数据集市B — 数据集市C\par

4. 企业级数据集市结构则强调从企业范围统一规划多个数据集市。\par

6. 数据仓库的间接访问\par

PPT 中强调，有时存在从数据仓库环境到操作型环境的数据回流，但操作型环境中采用 SQL 等方式直接访问数据仓库的“回流”是不正常的。直接访问有响应时间、数据量、技术一致性、格式化等限制，因此直接访问很少。\par

\newpage
\section*{12、间接访问的基本思想}
\addcontentsline{toc}{section}{12、间接访问的基本思想}

数据仓库的间接访问不是让操作型系统实时查询庞大的数据仓库，而是：\par

后台程序周期性分析数据仓库；\par
生成一个较小、易访问的在线预分析数据文件；\par
操作型环境访问这个小文件，而不是直接访问整个数据仓库。\par

两个例子：\par

航空公司佣金计算系统：\par
最佳佣金计算需要当前预定情况和历史订票情况，航空公司离线完成佣金计算和航班历史分析，并创建小型航班状态表，供在线交互时快速查询。\par

零售个性化系统：\par
销售代表与顾客电话咨询时，需要立即查询顾客个性化信息，如上次购物日期、上次购物类型、市场类别信息等。这些信息由数据仓库后台分析程序定时提供。\par

在线预分析文件特点\par

每个数据单元只包含少量数据\par
总体可包含大量信息\par
准确包含在线处理人员需要的内容\par
不被修改\par
批量方式周期性更新\par
是在线高性能访问环境的一部分\par
访问效率高\par
适合访问单个数据单元，而不是大量数据\par

\newpage
\section*{13、OLAP的评价准则}
\addcontentsline{toc}{section}{13、OLAP的评价准则}

Codd 准则：\par
1. 多维概念视图\par
OLAP 必须能从多个维度考察对象。\par
2. 透明性准则\par
OLAP 在体系结构中的位置和数据源对用户透明。\par
3. 存取能力准则\par
能把 OLAP 概念视图映射到异质数据存储上，并提供单一、完整、连续的用户视图。\par
4. 稳定的报表功能\par
当维数和综合层次增加时，报表能力和响应速度不应明显下降。\par
5. 客户 / 服务器体系结构\par
服务器负责统一概念模式、逻辑模式和物理模式；客户端负责应用逻辑和界面。\par
6. 维的同等性原则\par
每一数据维在数据结构和操作能力上都是等同的。\par
7. 动态稀疏矩阵处理准则\par
OLAP 工具应适应特定模型的数据分布，尤其要处理稀疏数据。\par
8. 多用户支持能力准则\par
支持并发访问、数据完整性和安全性机制。\par
9. 非受限的跨维操作\par
对维之间固有层次关系可以自动推导；其他计算需要提供计算完备的语言。\par
10. 直观的数据操作\par
11. 灵活的报表生成\par
12. 不受限维与聚集层次\par
维数不应小于 15，并支持任意聚集层次。\par

FASMI 准则\par

Fast：快速\par
Analysis：分析\par
Shared：共享\par
Multidimensional：多维\par
Information：信息\par

\newpage
\section*{14、OLAP 的基本数据模型}
\addcontentsline{toc}{section}{14、OLAP 的基本数据模型}

主要分为两类：\par
MOLAP：多维联机分析处理\par
ROLAP：关系联机分析处理\par

1. MOLAP：多维数据库模型\par

MOLAP 全称是 Multi-Dimensional OLAP，即多维联机分析处理。\par
MOLAP 使用专门的多维数据库管理系统管理数据。多维数据库采用的基本数据模式是多维数组。\par

例如一个销售分析可以表示为：\par

（时间，商品，商店，销售额）\par

其中：\par

时间、商品、商店是维度；\par
销售额是度量值；\par
每一个维度取一个值，就可以定位到一个具体的数据单元。\par

例如：\par

（2024年1月，电视机，南京店，50000元）\par

这就是一个多维数据单元。\par

MOLAP 使用专门的多维数据库存储数据，将事实表表示成多维数组，把维属性值映射为数组下标，把度量值存储在数据单元中。其优点是查询和综合速度快，缺点是预计算可能造成数据爆炸，维数和动态扩展能力受到限制。\par

2. ROLAP：关系数据库模型\par

ROLAP 全称是 Relational OLAP，即关系联机分析处理。\par

ROLAP 用传统的关系数据库管理系统 RDBMS 管理数据，将星型模型或雪花模型用二维表形式存储，表之间通过关键字连接，形成一个关系模式。ROLAP 不专门使用多维数据库，而是用关系数据库中的表来模拟多维结构。\par

例如销售分析可以设计为：\par

销售事实表 + 时间维表 + 产品维表 + 商店维表\par

用户在 ROLAP 上的查询操作，会被改写成关系数据库中的 SQL 查询来执行。\par

ROLAP 使用关系数据库管理 OLAP 数据，把星型或雪花模式用二维关系表存储，表之间通过关键字连接。用户的多维查询会被转换为关系数据库查询执行。其优点是适合大规模数据和维度较多的情况，缺点是查询性能通常不如 MOLAP。\par

\newpage
\section*{14、星型模型和雪花模型}
\addcontentsline{toc}{section}{14、星型模型和雪花模型}

一、星型模型\par

星型模型是 ROLAP 中最典型、最常用的数据模型。\par

星型模式是一种多维表结构，它一般由事实表和维表组成。事实表存放多维表中的主要事实，也就是度量值；维表存放维成员的取值。一个 n 维的多维表通常有 n 个维表和一个事实表，它们构成一个星形结构，所以称为星型模式。\par

1. 星型模型的组成\par

星型模型由两类表组成：\par

\begin{tabularx}{\linewidth}{|*{2}{>{\RaggedRight\arraybackslash}X|}}
\hline
表类型 & 作用 \\
\hline
事实表 Fact Table & 存放度量值，以及连接各维表的外键 \\
\hline
维表 Dimension Table & 存放维度成员及其描述属性 \\
\hline
\end{tabularx}

2. 事实表\par

事实表是星型模型的中心。\par

它通常包括：\par

各个维度的外键；\par
一个或多个度量值。\par

例如销售表可以表示为：\par

销售表（产品标识符，商店标识符，日期标识符，销售额）\par

其中：\par

产品标识符、商店标识符、日期标识符是外键；\par
销售额是度量值。\par

3. 维表\par

维表围绕事实表展开，用来描述分析角度。\par

产品表（产品标识符，类别，大类别）\par
商店表（商店标识符，市名，省名，国名，洲名）\par
时间表（时间标识符，日期，月份，季度，年份）\par

这些维表分别提供：\par

产品维度\par
地域 / 商店维度\par
时间维度\par

4. 星型模型结构示意\par

可以这样理解：\par

\hspace*{1em}产品维表\par
\hspace*{1em}|\par
\hspace*{1em}|\par
时间维表 —— 销售事实表 —— 商店维表\par

中心是事实表，周围是维表，结构像星星，所以叫星型模型。\par

5. 星型模型的特点\par

\begin{tabularx}{\linewidth}{|*{2}{>{\RaggedRight\arraybackslash}X|}}
\hline
特点 & 说明 \\
\hline
结构简单 & 一个事实表直接连接多个维表 \\
\hline
查询效率较高 & 连接路径短，适合 OLAP 查询 \\
\hline
面向分析 & 用户可以从不同维度观察事实 \\
\hline
维表通常反规范化 & 维表中会保留层次属性，存在一定冗余 \\
\hline
易于理解 & 结构直观，适合报表和多维分析 \\
\hline
\end{tabularx}

星型模型由一个中心事实表和多个维表组成。事实表保存度量值和维度外键，维表保存维成员及其描述属性。事实表与各维表通过公共关键字连接，形成星形结构。星型模型结构简单、查询连接少、分析效率高，是 ROLAP 中常用的数据模型。\par

三、雪花模型\par

雪花模型是星型模型的扩展。\par

雪花模型对星型模型的维表进一步层次化，原有的各维表可能被扩展为小的事实表或多个层次表，形成一些局部的“层次”区域。\par

简单说，雪花模型就是把星型模型中的维表继续拆分，使维表更加规范化。\par

1. 雪花模型的形成\par

在星型模型中，产品维表可能是：\par

产品表（产品标识符，类别，大类别）\par

如果进一步规范化，可以拆成：\par

产品表（产品标识符，类别编号）\par
类别表（类别编号，大类别编号）\par
大类别表（大类别编号，大类别名称）\par

这样产品维度就不再是一张宽表，而是被拆分成多个有层次关系的小表。\par

2. 雪花模型例子\par

销售事实表（产品no，商店no，日期no，销售额）\par

产品（产品no，类别no，供应商no，顾客no）\par
类别（类别no，层1，层2，……）\par
顾客（顾客no，年龄层次，职业类型）\par
年龄层次（名称）\par
职业类型（名称）\par
供应商（供应商no，供应商类别，……）\par

可以看到，产品维表被拆成类别、顾客、供应商等更细的维度层次表。\par

3. 雪花模型结构示意\par
\hspace*{1em}类别表\par
\hspace*{1em}|\par
产品维主表 —— 销售事实表 —— 商店维表\par
\hspace*{1em}|\par
\hspace*{1em}供应商表\par
\hspace*{1em}|\par
\hspace*{1em}供应商类别表\par

整体结构不像星星那么简单，而是分叉更多，像雪花一样。\par

4. 雪花模型的优点\par

雪花模型的优点是：\par

使维表尽可能规范化，最大限度地减少维表的数据存储量。\par

也就是\par

减少数据冗余；\par
节省存储空间；\par
层次结构更清楚；\par
适合维度层次复杂的场景。\par

5. 雪花模型的缺点\par

执行查询时需要更多连接操作，可能影响查询性能。\par

也就是\par

表更多；\par
Join 更多；\par
查询路径更长；\par
用户理解成本更高；\par
OLAP 查询效率可能低于星型模型。\par

雪花模型是星型模型的扩展，它将星型模型中的维表进一步层次化和规范化，把一个维表拆分为多个具有层次关系的小表。其优点是减少维表冗余、节省存储空间；缺点是查询时连接操作增加，可能降低查询性能。\par

\newpage
\section*{15、数据立方体}
\addcontentsline{toc}{section}{15、数据立方体}

数据立方体来源于物化视图思想。PPT 中说，数据仓库的数据模式可以看成定义在多个数据源上的数据视图，其中分析型数据通常是统计得到的综合数据。如果每次都用统计函数临时计算，时间开销太大，因此可以预先计算统计信息并存入数据仓库，这就是物化视图；存放物化视图的三维数据模型叫数据立方体。\par

基本理解\par

以销售事实表为例，可以有三个维度：\par

产品 Product\par
商店 Store\par
日期 Date\par

可以构造：\par

PSD 视图\par
PS、PD、SD 视图\par
P、S、D 视图\par
ALL 视图\par

其中 ALL 视图表示不分组，统计全部销售总额。\par

数据超立方体\par

当分析对象超过三维时，就构成 n 维应用，无法用普通三维立方体表示，只能通过虚拟 n 维空间建立 n 维立方体，即数据超立方体。\par

方体格\par

n 维方体称为基本方体\par
0 维方体代表最高层次抽象，称为顶点方体\par
方体格构成数据立方体\par
10.4 数据立方体的物化\par

几种物化方式：\par

不进行物化\par
全物化\par
部分物化\par
冰山方体\par
利用查询处理时物化的方体\par
在装入和刷新时有效更新物化方体\par

数据立方体是为了提高综合统计查询速度而预先计算并保存多维统计结果的数据结构，本质上是物化视图在多维数据分析中的组织形式。它可以包含基本方体、低维方体和顶点方体，并可采用不物化、全物化、部分物化和冰山方体等方式进行存储优化。\par

\newpage
\section*{16、多维数据分析}
\addcontentsline{toc}{section}{16、多维数据分析}

多维分析是对以多维形式组织起来的数据进行切片、切块、旋转、钻取等分析动作，使用户能从多个角度、多个侧面观察数据，从而深入理解数据中的信息。\par

基本概念\par
度量值：分析型处理中关心和分析的对象，通常是数字值，如销售金额、成本金额、库存数量、消费金额。\par

维：观察度量值的角度，如时间维、商品维、地域维。\par

层：维中的分析深度。例如时间维可以有日、月、季、年；地域维可以有市、省、国、洲。\par

维成员：维的一个取值。如时间维中的“1998 年”“1998 年 1 月”，地域维中的“江苏”“南京”。\par

数据单元：当多维数组的每一维都选定一个维成员时，就唯一确定一个度量值，这个值或存放该值的位置称为数据单元。\par

基本操作\par

\begin{tabularx}{\linewidth}{|*{2}{>{\RaggedRight\arraybackslash}X|}}
\hline
操作 & 含义 \\
\hline
切片 Slice & 根据某一维上的某个维成员值选择统计数据 \\
\hline
切块 Dice & 根据一个或多个维度上的维成员取值区间选择统计数据 \\
\hline
旋转 Pivot / Rotate & 调整维度排列次序，改变观察角度 \\
\hline
上钻 / 数据概括 Roll-up & 把多维下标提升到较高概念层次，形成更概括的统计结果 \\
\hline
下钻 / 数据细化 Drill-down & 把多维下标降低到较低概念层次，形成更细致的统计结果 \\
\hline
跨钻 Drill across & 对多个事实进行操作 \\
\hline
钻透 Drill through & 下钻到数据立方体最低细节后，继续细化到数据仓库 / 数据库关系表 \\
\hline
\end{tabularx}

例子：\par
对销售额进行分析时，可以从“年”下钻到“季度”“月”“日”；也可以从“商品大类”下钻到“小类”“具体商品”；还可以把报表中的时间维和地区维调换位置，这就是旋转。\par

\newpage
\section*{17、数据仓库设计的原则}
\addcontentsline{toc}{section}{17、数据仓库设计的原则}

面向主题原则\par
数据驱动原则\par
原型法设计原则\par

数据仓库设计的三条原则是面向主题、数据驱动和原型法。面向主题强调从决策分析主题出发组织数据；数据驱动强调从已有操作型数据源出发进行抽取、综合和集成；原型法强调在需求不明确且不断变化的情况下，先建立基本框架，再通过用户反馈迭代完善。\par

\newpage
\section*{18、数据仓库设计的三级数据模型}
\addcontentsline{toc}{section}{18、数据仓库设计的三级数据模型}

概念模型、逻辑模型、物理模型\par

1. 概念模型\par

概念模型是为一定目标设计系统、收集信息而服务的概念型工具，是客观世界到机器世界的中间层次。PPT 中提到可以使用 E-R 法表示概念模型。\par

概念模型设计主要做三件事：\par

\begin{tabularx}{\linewidth}{|*{2}{>{\RaggedRight\arraybackslash}X|}}
\hline
内容 & 说明 \\
\hline
确定系统边界 & 明确要支持哪些决策、决策者关心什么问题、需要什么信息、涉及哪些数据源 \\
\hline
确定主要主题及其内容 & 明确数据仓库的分析对象，即主题 \\
\hline
OLAP 等分析应用设计 & 明确后续要支持哪些分析操作 \\
\hline
\end{tabularx}

确定主题时，需要描述主题及其属性信息、属性取值情况、主题公共码键、主题之间的关系，之后可以用 E-R 模型表示概念数据模型。\par

例子：\par
商品、顾客、供应商可以作为主题。商品主题可以包含商品固有信息、采购信息、销售信息、库存信息等。\par

2. 逻辑模型\par

逻辑模型描述数据仓库主题的逻辑实现，PPT 中对应的是关系模型。\par

逻辑模型设计主要包括：\par

\begin{tabularx}{\linewidth}{|*{2}{>{\RaggedRight\arraybackslash}X|}}
\hline
内容 & 说明 \\
\hline
E-R 图转换为关系表 & 把概念模型转换成二维表 \\
\hline
定义数据源和抽取规则 & 明确每个属性来自哪个源系统、源表、源字段 \\
\hline
粒度划分 & 确定数据是细节数据、轻度总结还是高度总结 \\
\hline
数据分割 & 把逻辑整体数据分割成较小、可独立管理的数据单元 \\
\hline
定义合适的数据来源 & 确保数据可获得、可集成、可分析 \\
\hline
\end{tabularx}

粒度划分会直接影响数据仓库的数据量和能够回答的查询类型，是影响数据仓库性能的重要因素。\par

3. 物理模型\par

物理模型是逻辑模型在数据仓库中的实现。PPT 中说，物理模型设计是在逻辑模型基础上确定数据的存储结构、索引策略、存储分配、数据存放位置等物理内容，目的是提高数据仓库访问性能。\par

常用技术包括：\par

\begin{tabularx}{\linewidth}{|*{2}{>{\RaggedRight\arraybackslash}X|}}
\hline
技术 & 作用 \\
\hline
合并表 & 减少多表连接的 I/O 开销 \\
\hline
建立数据序列 & 提高顺序访问效率 \\
\hline
引入冗余 & 减少分析时需要访问的表数量 \\
\hline
表的物理分割 & 提高管理和查询性能 \\
\hline
生成导出数据 & 预先生成常用分析数据 \\
\hline
建立广义索引 & 加快查询 \\
\hline
规范化 / 反规范化 & 在存储与查询性能之间权衡 \\
\hline
\end{tabularx}

数据仓库设计的三级数据模型包括概念模型、逻辑模型和物理模型。概念模型是客观世界到机器世界的中间层次，用于确定系统边界、主要主题及分析应用，常用 E-R 法表示；逻辑模型描述数据仓库主题的逻辑实现，通常把 E-R 图转换为关系表，并定义数据源、抽取规则、粒度和数据分割策略；物理模型是逻辑模型在数据仓库中的具体实现，关注存储结构、索引策略、存储分配和物理优化，其目的是提高数据仓库访问性能。\par

\newpage
\section*{19、数据挖掘的方法}
\addcontentsline{toc}{section}{19、数据挖掘的方法}

定义：对数据库或数据仓库中蕴涵的、未知的、非平凡的、有潜在应用价值的模式或规则的提取\par

\begin{tabularx}{\linewidth}{|*{2}{>{\RaggedRight\arraybackslash}X|}}
\hline
方法 & 作用 \\
\hline
特征规则挖掘 & 描述一类数据对象的普遍特征 \\
\hline
关联规则挖掘 & 发现事务数据中项与项之间的关联关系 \\
\hline
序列模式分析 & 发现数据对象之间的前后时序关系 \\
\hline
分类分析 & 根据训练样本建立模型，对未知数据分类 \\
\hline
聚类分析 & 对没有类别标记的数据自动划分群组 \\
\hline
\end{tabularx}

1. 特征规则挖掘\par
是一种常见的知识形式，它用于描述一类数据对象的普遍特征，是普化知识的一种\par

特征规则的数据挖掘方法有两类：\par

1. 面向属性归约方法\par

它通过属性值之间的概念层次结构进行归约，获得概括性知识。例如，把“南京”提升为“江苏”，再把“江苏”提升为“华东”，这叫概念提升。概念提升的目的包括规范化属性取值，提高模式的置信度和兴趣度。\par

2. 数据立方方法\par

数据立方方法利用预先计算好的统计、求和等综合结果进行挖掘，可以减少重复统计查询，提高挖掘效率。它和 OLAP 中的上卷、下钻思想有关。\par

特征规则挖掘用于描述一类对象的普遍特征，属于普化知识发现。其主要方法包括面向属性归约方法和数据立方方法。面向属性归约通过概念层次树对属性值进行概念提升，以获得概括性知识。\par

2. 关联规则挖掘\par

关联规则挖掘用于发现事务数据库中多个属性或项之间的关联程度。\par

典型形式是：\par

A → B\par

含义是：出现 A 的事务中，经常也出现 B。\par

例如：\par

购买尿布 → 购买啤酒\par
购买面包和黄油 → 购买牛奶\par

关联规则挖掘用于发现事务数据中项与项之间的关联关系，规则形式为 A→B。其核心指标是支持度和置信度，满足最小支持度和最小置信度的规则称为强规则。Apriori 算法利用“频繁项集的任一子集必定频繁”的性质寻找频繁项集。\par

3. 序列模式分析\par

序列模式分析与关联规则类似，也是为了找出数据对象之间的联系，但它更强调前后时序关系。也就是说，它不仅关心“哪些项经常一起出现”，还关心“哪些事件先发生、哪些事件后发生”。\par

例如：\par

下雨 → 洪涝\par
购买电筒 → 购买电池\par

序列模式分析用于发现具有时间先后关系的数据对象之间的联系。它与关联规则挖掘类似，但更强调事件之间的前因后果和时序关系。\par

4. 分类分析\par

分类分析是数据挖掘的重要方法之一。分类分析通过分析训练数据样本，产生关于类别的精确描述，并用于对未来数据进行分类和预测。\par

1. 基本思想\par

分类分析需要先有已经标记类别的数据。例如：\par

信用等级：优、良、一般、差；\par
是否购买电脑：yes / no；\par
客户是否流失：流失 / 未流失。\par

系统先根据这些带类别标记的数据学习规则，再对未知类别的数据进行预测。\par

2. 分类分析的两个步骤\par
\begin{tabularx}{\linewidth}{|*{2}{>{\RaggedRight\arraybackslash}X|}}
\hline
步骤 & 内容 \\
\hline
建立模型 & 用训练数据集构造分类模型 \\
\hline
使用模型分类 & 评估模型准确性，并对未知类别数据分类 \\
\hline
\end{tabularx}

因为训练样本已经给出了类标号，所以分类分析属于有指导学习。\par

3. 常见分类方法\par

包括：\par
决策树方法\par
贝叶斯方法\par
人工神经网络方法\par
约略集方法\par
遗传算法\par

其中，决策树是一种树结构：内部结点表示属性测试，边表示测试结果，叶结点表示类别或类别分布。\par

分类分析是一种有指导学习方法。它先利用带有类标号的训练数据建立分类模型，再用该模型对未知类别的数据进行分类和预测。常见方法包括决策树、贝叶斯、人工神经网络、约略集和遗传算法。\par

5. 聚类分析\par

聚类分析又称集群分析。它和分类分析相反：分类分析的数据已经有类别标记，而聚类分析面对的是没有类别标记的记录。系统按照距离或相似性等规则自动划分记录集合，相当于自动给数据打标签。\par

聚类分析可以分为距离聚类和相似系数聚类，主要方法包括：\par

划分方法\par
层次方法\par
基于密度的方法\par
基于网格的方法\par
基于模型的方法\par

聚类分析的结果可以表现为聚类树\par

聚类分析是一种无指导学习方法，输入数据没有预先给定类别标记，系统根据距离或相似性自动把记录划分为若干簇。主要方法包括划分方法、层次方法、基于密度的方法、基于网格的方法和基于模型的方法。\par

总结：常用的数据挖掘方法包括特征规则挖掘、关联规则挖掘、序列模式分析、分类分析和聚类分析等。特征规则挖掘用于描述一类数据对象的普遍特征，常用方法有面向属性归约方法和数据立方方法；关联规则挖掘用于发现事务数据中项与项之间的关联关系，规则形式为 A→B，核心指标是支持度和置信度；序列模式分析用于发现数据对象之间的前后时序关系；分类分析是一种有指导学习方法，通过带类标号的训练数据建立模型，并对未知类别数据进行分类和预测；聚类分析是一种无指导学习方法，对没有类别标记的数据按照距离或相似性自动划分为若干簇。此外，数据挖掘还包括异常分析、时序数据挖掘等方法。\par

\newpage
\section*{19、多维建模}
\addcontentsline{toc}{section}{19、多维建模}

多维建模的基本概念包括：事实表、维度表、事实与维度的融合、星型模型、雪花模型、数据立方体。\par

1. 事实表\par

事实表是维度建模的核心和基本表。每一事实表都对应一个或若干个度量值，度量值是事实表的核心，也是趋势分析的对象。事实表用来记录维度值与度量值之间的关系。\par

事实表的特点：\par

\begin{tabularx}{\linewidth}{|*{2}{>{\RaggedRight\arraybackslash}X|}}
\hline
特点 & 说明 \\
\hline
核心表 & 维度建模的中心 \\
\hline
保存度量值 & 如销售量、销售额、成本额、毛利润 \\
\hline
保存外键 & 通过外键连接维表 \\
\hline
粒度一致 & 事实表中的所有度量值必须具有相同粒度 \\
\hline
主键 & 通常可由多个外键组合构成 \\
\hline
\end{tabularx}

每个事实表都有两个或两个以上外关键字，通过外关键字建立事实表与维表之间的联系。\par

事实表的粒度模型\par

事实表粒度划分模型包括：\par

\begin{tabularx}{\linewidth}{|*{2}{>{\RaggedRight\arraybackslash}X|}}
\hline
粒度模型 & 说明 \\
\hline
事务事实表 & 每个业务事务一行 \\
\hline
周期快照事实表 & 按固定周期记录状态 \\
\hline
累积快照事实表 & 跟踪一个过程从开始到结束的多个阶段 \\
\hline
\end{tabularx}

库存管理课件中也具体举例：库存事实表可以有库存周期快照、库存事务、库存累积快照三种互补模型。\par

2. 度量值\par

事实表中的度量值通常是数值类型，是可以连续取值的量，很少采用文本形式。度量值分为三类：\par

\begin{tabularx}{\linewidth}{|*{3}{>{\RaggedRight\arraybackslash}X|}}
\hline
类型 & 含义 & 例子 \\
\hline
可做加法运算 & 可以沿所有维度汇总 & 销售额、销售量 \\
\hline
可沿某些维度做加法运算 & 只能在部分维度上加总 & 库存量、账户余额 \\
\hline
不能做加法运算 & 不能直接求和 & 单价、毛利润率 \\
\hline
\end{tabularx}


3. 维度表\par

维度表是事实表的入口，为用户提供使用数据仓库的接口。维度表中的属性通常用于定义事实表上的查询条件，也可以作为报表和统计查询的列。\par

维度表特点：\par

\begin{tabularx}{\linewidth}{|*{2}{>{\RaggedRight\arraybackslash}X|}}
\hline
特点 & 说明 \\
\hline
事实表入口 & 用户通过维度属性访问事实表 \\
\hline
描述性强 & 保存文本型、离散型属性 \\
\hline
列多行少 & 相对于事实表，维表通常列多、行少 \\
\hline
用于查询条件 & 如时间、产品、商场、促销等 \\
\hline
\end{tabularx}

度量值属性有许多取值可能，并参与统计运算；维度属性通常是离散的、取值可能不多、很少变化、不参与统计计算但常用作查询条件。\par

4. 事实与维度的融合\par

多维建模的核心结构是：事实表通过关键字与相关维表连接。以日销售情况事实表为例：事实表中包含日期关键字、产品关键字、商场关键字等外键，并通过这些外键连接日期维度、产品维度和商场维度。\par

可以理解为：\par

日期维度\par
\hspace*{1em}|\par
产品维度 —— 销售事实表 —— 商场维度\par

用户通过维表中的属性过滤、分组和观察事实表中的度量值。\par

5. 多维建模的设计过程\par

\begin{tabularx}{\linewidth}{|*{2}{>{\RaggedRight\arraybackslash}X|}}
\hline
步骤 & 内容 \\
\hline
第一步 & 选取要建模的业务处理过程 \\
\hline
第二步 & 定义业务处理的粒度 \\
\hline
第三步 & 选择事实表中的维度 \\
\hline
第四步 & 选择事实表中的度量值 \\
\hline
\end{tabularx}

在零售营销案例中，用 POS 事务说明：\par

业务处理：商品销售；\par
粒度：POS 事务的单个商品条目；\par
维度：日期、商场、产品、促销；\par
事实：销售量、销售额、成本额、毛利润等。\par

6. 典型维度设计\par

日期维度是每个数据仓库必须具备的维度。\par

日期维度表可以事先建立，可以预先建立 5 到 10 年的日期维度值。\par

常见维度包括：\par

\begin{tabularx}{\linewidth}{|*{2}{>{\RaggedRight\arraybackslash}X|}}
\hline
维度 & 典型属性 \\
\hline
日期维度 & 日、周、月、季度、年、财政月、财政季度、节假日标志 \\
\hline
产品维度 & 产品描述、SKU 编号、小类、大类、部门、包装类型 \\
\hline
商场维度 & 商场名称、编号、地址、地区、销售面积 \\
\hline
促销维度 & 促销名称、促销媒体类型、促销起止日期 \\
\hline
顾客维度 & 顾客姓名、地址、分类属性等 \\
\hline
\end{tabularx}

7. 退化维度\par

\hspace*{1em}POS 事务编号时提出退化维度：维度表为空，具体维度值直接存放在事实表中。\par

例子包括：\par

事务编号\par
订单编号\par
发票编号\par
提货单编号\par

零售模型中，POS 事务编号就是退化维度。\par

多维建模是数据仓库中面向分析的数据组织方法，其基本概念包括事实表、维度表、星型模型、雪花模型和数据立方体。事实表是维度建模的核心，保存度量值和连接维表的外键，事实表中的所有度量值必须具有相同粒度；维度表是事实表的入口，保存描述性属性，用于查询条件、报表列和分组分析。多维建模的设计过程包括选择业务处理过程、定义粒度、选择维度和选择度量值。典型事实表粒度模型包括事务、周期快照和累积快照。\par
\newpage
\end{document}
